%%%%%%%%%%%%%%%%%%%%%%%%%%%%%%%%%%%%%%%%%%%%%%%%%%%%%%%
% % Lines starting with % are comments, which are ignored.
% % This is a handy way of indicating the date and version of
% % your document, to wit:
% %
% % LaTeX sample file
% % Modified March, 2002
% %
%%%%%%%%%%%%%%%%%%%%%%%%%%%%%%%%%%%%%%%%%%%%%%%%%%%%%%%
\DocumentMetadata{
  % PDF standards: ONE key whose value is a comma-separated list --
  % UA-2 (accessibility) plus A-4f (archival; the archival companion of
  % UA-2). Do NOT use separate pdfstandard=... keys: with ordinary
  % key-value parsing only the last one takes effect.
  pdfstandard={UA-2,A-4f},
  lang=en-US,
  tagging=on,
  tagging-setup={
    math/setup=mathml-SE
  }
}

\documentclass[12pt]{exam}
% worksheet-macros.tex -- shared package setup for worksheet-template.tex.
%
% NOTE ON PATHS: every template in this repository is stored one level up
% from where it is compiled (see README.md). All relative \input paths
% below are resolved against the folder the document is compiled in,
% which sits one directory above this file.

% --- Math ------------------------------------------------------------
% amsmath: equation environments and numbering
% amsfonts: \mathbb etc.  |  mathrsfs: \mathfrak  |  stmaryrd: extra symbols
\usepackage{amsmath}
\usepackage{amsfonts}
\usepackage{mathrsfs}
\usepackage{stmaryrd}

% --- Layout and glyphs ------------------------------------------------
% float: [H] float placement (stop floats where you put them)
% fontspec: Unicode/OpenType fonts (compile with XeLaTeX or LuaLaTeX)
% emoji: emoji glyphs in the text
\usepackage{float}
\usepackage{fontspec}
\usepackage{emoji}

% --- Typography --------------------------------------------------------
% Atkinson Hyperlegible as the document-wide sans serif (sfdefault):
% a high-legibility face that is easier to read for students with
% dyslexia or low vision.
\usepackage[sfdefault]{atkinson}

% --- Colors -------------------------------------------------------------
% dvipsnames: extended named color palette (the Accessible* colors in
%   lecture-macros.tex are built on it);
% table: \rowcolors / \cellcolor for shading worksheet tables.
% NOTE: xcolor may only be loaded ONCE per document -- a second
% \usepackage{xcolor} with different options is a fatal "Option clash"
% error, so ALL of its options are listed here.
\usepackage[dvipsnames,table]{xcolor}

% --- Multimedia ---------------------------------------------------------
% animate: embeds audio/video (and legacy Flash) media in the worksheet.
\usepackage{animate}

% --- Bibliography --------------------------------------------------------
% biblatex (biber backend): citations and a reference list for worksheets.
\usepackage[
  backend=biber,
  style=alphabetic,
  url=true,
  doi=true,
  isbn=true,
  backref=true
]{biblatex}

% --- Links ---------------------------------------------------------------
% hyperref makes every \ref/citation/URL clickable in the PDF.
% Link colors must satisfy WCAG 1.4.3 (>= 4.5:1 contrast on white):
%   blue = 8.6:1, green!40!black = 7.2:1.
% (The old cyan/magenta values measured 1.25:1 / 3.1:1 and failed.)
\usepackage{hyperref}
\hypersetup{
  colorlinks=true,
  linkcolor=blue,
  filecolor=blue,
  urlcolor=green!40!black,
}

% --- Shared course macros ------------------------------------------------
% Root macros.tex: math shorthands used across all course documents
% (\RR, \NN, \ZZ, \QQ, \grad, \curl, \tr, \der, \pder, ...).
\newcommand{\RR}{\mathbb{R}}
\newcommand{\NN}{\mathbb{N}}
\newcommand{\ZZ}{\mathbb{Z}}
\newcommand{\QQ}{\mathbb{Q}}
\newcommand{\vv}{\vec{v}}
\newcommand{\ww}{\vec{w}}
\newcommand{\xx}{\vec{x}}
\newcommand{\yy}{\vec{y}}
\newcommand{\zz}{\vec{z}}
\newcommand{\uu}{\vec{u}}
\newcommand{\basis}[1]{\vec{e_{#1}}}
\newcommand{\id}{\mathrm{id}}
\newcommand{\CC}{\mathbb{C}}
\renewcommand{\div}{\operatorname{div}}
\newcommand{\grad}{\operatorname{grad}}
\newcommand{\curl}{\operatorname{curl}}
\newcommand{\tr}{\operatorname{tr}}
\newcommand{\der}[2]{\dfrac{d #1}{d #2}}
\newcommand{\pder}[3][]{\dfrac{\partial^{#1} #2}{\partial #3}}


% % Title and author(s)
%%%%%%%%%%%%%%%%%%%%%%%%%%%%%%%%%%%%%%%%%%%%%%%%%%%%%%%
\title{Worksheet 13: Related Rates}
\author{MATH 1271}
\date{10/22/2026}
%%%%%%%%%%%%%%%%%%%%%%%%%%%%%%%%%%%%%%%%%%%%%%%%%%%%%%%
% %
% % The next command allows your in import encapsulated
% % postscript files, .epsf or .eps files, which
% % contain vector graphic image data.
% %
%%%%%%%%%%%%%%%%%%%%%%%%%%%%%%%%%%%%%%%%%%%%%%%%%%%%%%%
% %   The first thanks indicates your affiliation
% %
% %  Just the name here.
% %
% % Your mailing address goes at the end.
% %
% % \thanks is also how you indicate grant support
% %
%%%%%%%%%%%%%%%%%%%%%%%%%%%%%%%%%%%%%%%%%%%%%%%%%%%%%%%

%%%%%%%%%%%%%%%%%%%%%%%%%%%%%%%%%%%%%%%%%%%%%%%%%%%%%%%
% $ The next command ensures the margins are 1 inch on
% % all sides.
% %
%%%%%%%%%%%%%%%%%%%%%%%%%%%%%%%%%%%%%%%%%%%%%%%%%%%%%%%
\begin{document}
\newpage
\maketitle

\begin{enumerate}

\item A water tank has the shape of an inverted circular cone (point down, as in an example from lecture) with a base of radius 6 feet and a depth (``height") of 8 feet. Suppose that water is being pumped into the tank at a constant instantaneous rate of 4 cubic feet per minute. Let $r$ and $h$ be the radius and the depth of the water at a given time, $t$.

\begin{enumerate}

\item Draw a picture of the conical tank, including a sketch of the water level at a point in time when the tank is not yet full. Indicate what the variables $r$ and $h$ represent on your figure.

\vfill

\item What equation relates the radius and height of the water? (Hint: Similar triangles!)

\vfill

\item Determine an equation that relates the volume of water, $V$, in the tank at time $t$ to the depth of the water, $h$, at that time. 

\vfill

\newpage

\item Through differentiation, find an equation that relates the instantaneous rate of change of the water volume, $V$ (with respect to time) to the instantaneous rate of change of the depth of the water.

\vfill

\item Find the rate at which the water level is rising at the moment when the water in the tank is $3$ feet deep. 

\vfill

\item When is the water rising most rapidly? When the depth is $h=3$, $h=4$, or $h=5$ feet?

\vfill

\end{enumerate}

\newpage




\item A television camera is positioned 4000 feet from the base of a rocket launching pad. The angle of elevation of the camera has to change at the correct rate in order to keep the rocket in sight. In addition, the auto-focus of the camera has to take into account the increasing distance between the camera and the rocket. We assume that the rocket rises vertically. 

(A similar problem is discussed and pictured dynamically \url{https://gvsu.edu/s/9t/}. Exploring the applet at the link will be helpful to you in answering the questions that follow.)

\begin{enumerate}

\item Draw a figure that summarizes the given situation. What parts of the picture are changing? What parts are constant? Introduce variables to represent the quantities that are changing.

\vfill

\item Find an equation that relates the camera’s angle of elevation to the height of the rocket.

\vfill

\item Find an equation that relates the instantaneous rate of change of the camera’s elevation angle to the instantaneous rate of change of the rocket’s height (where all rates of change are with respect to time).

\vfill

\newpage


\item Find an equation that relates the rocket’s height to the distance between the camera and the rocket.

\vfill

\item Find an equation that relates the instantaneous rates of change of the rocket’s height to the instantaneous rate of change of the distance between the camera and the rocket.

\vfill


\newpage

\item Suppose that the rocket’s speed is 600 ft/sec at the instant it has risen 3000 feet. 

\begin{enumerate}

\item How fast is the distance from the television camera to the rocket changing at that moment? 

\vfill

\item If the camera is following the rocket, how fast is the camera’s angle of elevation changing at that same moment?

\vfill

\end{enumerate}

\item If from an elevation of 3000 feet onward the rocket continues to rise at 600 feet/sec, will the rate of change of distance with respect to time be greater when the elevation is 4000 feet than it was at 3000 feet, or less? Why?

\vfill

\end{enumerate}

\newpage

\item Suppose you are driving along some interstate where the speed limit is 70 miles per hour. A state trooper is hidden about $\dfrac{3}{100}$'ths of a mile away from the highway, and uses a radar gun to measure the rate of change of a vehicle's distance away from them.

\begin{enumerate}
  \item Draw a picture of this situation and label any quantities you think might be important to this situation.
    \vfill

  \item Suppose that you are a distance of $\dfrac{5}{100}$'ths of a mile away from the state trooper,
    determine the maximum speed at which you can drive so that the state trooper
    would still measure your speed below the speed limit.
    \vfill
\end{enumerate}

\end{enumerate}

% Most people don't use these, so they are "commented out"
% by starting the lines with a "%"
%\begin{keywords}
%   \LaTeX, typesetting
%\end{keywords}

%\begin{AMS}
%   50C60, 18C25
%\end{AMS}

%%%%%%%%%%%%%%%%%%%%%%
% % Here is the start of the Text
%%%%%%%%%%%%%%%%%%%%%%
\end{document}
